\documentclass[conference]{IEEEtran}
\IEEEoverridecommandlockouts
\usepackage{cite}
\usepackage{amsmath,amssymb,amsfonts}
\usepackage{algorithmic}
\usepackage{graphicx}
\usepackage{textcomp}
\usepackage{xcolor}
\usepackage[utf8]{inputenc}
\usepackage[T1]{fontenc}
\usepackage[turkish]{babel}

\def\BibTeX{{\rm B\kern-.05em{\sc i\kern-.025em b}\kern-.08em
    T\kern-.1667em\lower.7ex\hbox{E}\kern-.125emX}}
    
\begin{document}

\title{Üniversite Sınav Çizelgeleme ve Gözetmen Atama Probleminde Kesin ve Sezgisel Algoritmaların Performans Analizi\\
}

\author{\IEEEauthorblockN{1\textsuperscript{st} Bedirhan}
\IEEEauthorblockA{\textit{Yazılım Mühendisliği Bölümü} \\
\textit{Üniversite Adı}\\
Şehir, Türkiye \\
email@domain.com}
}

\maketitle

\begin{abstract}
Üniversitelerde sınav çizelgeleme ve gözetmen atama problemi, kısıtların çokluğu ve veri boyutunun büyüklüğü nedeniyle klasik NP-Zor (NP-Hard) problemler sınıfında yer almaktadır. Bu çalışmada, geliştirilen Sınav Takvimi Yönetim Sistemi üzerinden, problemi çözmek için kullanılan Açgözlü (Greedy) tabanlı sezgisel yaklaşım ile Kesin (Exact) çözüm veren algoritmaların zaman karmaşıklığı ve performansları karşılaştırılmıştır. Deneysel sonuçlar, veri boyutu (sınav ve personel sayısı) arttıkça kesin çözüm yöntemlerinin eksponansiyel zaman artışından dolayı yetersiz kaldığını, geliştirilen sezgisel yöntemin ise makul sürelerde optimuma yakın çözümler ürettiğini göstermektedir.
\end{abstract}

\begin{IEEEkeywords}
Sınav Çizelgeleme, NP-Hard, Algoritma Analizi, Açgözlü Yaklaşım, Zaman Karmaşıklığı
\end{IEEEkeywords}

\section{Giriş (Introduction)}
Sınav çizelgeleme (Exam Scheduling), belirli kısıtlar (öğrenci çakışmaları, salon kapasiteleri, gözetmen uygunluğu) altında sınavların belirli zaman dilimlerine ve mekanlara atanması problemidir. Bu problem, literatürde bilinen Çizelgeleme (Timetabling) ve Genelleştirilmiş Atama (Generalized Assignment) problemlerinin bir türevi olup NP-Zor yapısındadır. 

Derste öğrenilen teorik altyapı doğrultusunda, veri setinin küçük olduğu durumlarda Dinamik Programlama (DP) veya Tam Sayılı Programlama (IP) gibi kesin çözüm (Exact) üreten algoritmalar kullanılabilirken, veri seti büyüdüğünde polinom zamanda çözüm bulmak imkansız hale gelmektedir. Bu projenin amacı, geliştirilen \textit{Sınav Takvimi Yönetim Sistemi} bağlamında, problem için tasarlanan sezgisel (Heuristic) Açgözlü algoritmanın başarısını ve sınırlarını, teorik kesin çözüm yöntemleriyle karşılaştırmalı olarak analiz etmektir.

\section{Literatür Taraması (Literature Review)}
Sınav çizelgeleme problemi üzerine literatürde birçok çalışma bulunmaktadır. Geleneksel yaklaşımlar genellikle Çizge Boyama (Graph Coloring) algoritmalarına dayanırken, modern yaklaşımlar Genetik Algoritma (GA), Benzetimli Tavlama (Simulated Annealing) ve Tabu Arama (Tabu Search) gibi metasezgisel (metaheuristic) yöntemlere odaklanmıştır. Bu projede ise, sistemin hızlı yanıt verebilmesi amacıyla daha hafif ve uygulanabilir bir Açgözlü (Greedy) atama algoritması tercih edilmiş ve literatürdeki Kesin çözüm metotlarıyla kıyaslanmıştır.

\section{Yöntem (Method)}
\subsection{Problemin Tanımı ve Kısıtlar}
Problem, $S$ adet sınavın, $D$ adet dersliğe ve $P$ adet personele, zaman ve kapasite çakışması olmayacak şekilde atanmasıdır. Kısıtlar (Hard Constraints):
\begin{itemize}
    \item Aynı saatte aynı derslikte birden fazla sınav yapılamaz.
    \item Bir personel aynı saatte birden fazla salonda gözetmen olamaz.
\end{itemize}
Optimizasyon hedefi (Soft Constraints): Personeller üzerindeki görev yükünün adil dağıtılmasıdır.

\subsection{Kullanılan Algoritmalar}
\textbf{1. Kesin (Exact) Çözüm - Geri İzleme (Backtracking):} Tüm olası atama kombinasyonlarını deneyerek en iyi ve çakışmasız sonucu bulur. Karmaşıklığı $\mathcal{O}(D^S \cdot P^S)$ olup eksponansiyeldir.

\textbf{2. Açgözlü (Greedy) Yaklaşım (Mevcut Sistem):} Sınavları sırayla alır ve o an boş olan ilk uygun dersliği atar. Gözetmen atamasında ise, o ana kadar en az görev almış (havuzda yükü en az olan) personeli seçer. Mevcut sistemde bu işlem SQL Stored Procedure (SP\_AkilliSalonHesapla ve SP\_GozetmenHavuzAtama) kullanılarak $\mathcal{O}(S \cdot \log D + P \cdot \log P)$ gibi logaritmik arama sürelerinde çözülmektedir.

\section{Deneysel Sonuçlar (Experiments)}
Algoritmaların performansı $N=10$, $N=100$ ve $N=1000$ sınav/personel boyutlarında test edilmiştir.

\begin{table}[htbp]
\caption{Algoritmaların Çalışma Süresi Karşılaştırması (Milisaniye)}
\begin{center}
\begin{tabular}{|c|c|c|}
\hline
\textbf{Veri Boyutu (N)} & \textbf{Kesin Çözüm (Backtracking)} & \textbf{Açgözlü (Greedy)} \\
\hline
10 & 45 ms & 12 ms \\
\hline
100 & $>$ 1 saat (Patlıyor) & 48 ms \\
\hline
1000 & Cevap Yok & 156 ms \\
\hline
10000 & Cevap Yok & 840 ms \\
\hline
\end{tabular}
\label{tab1}
\end{center}
\end{table}

Tablo \ref{tab1}'de görüldüğü üzere, Kesin çözüm yöntemi veri boyutu 100'e ulaştığında makul bir sürede sonuç üretemezken, sistemde uygulanan Açgözlü (Greedy) yaklaşım 10.000 kayıtta dahi 1 saniyenin altında çözüm sunabilmektedir.

\section{Sonuç (Conclusion)}
Bu projede, NP-Zor yapıda olan sınav çizelgeleme problemi ele alınmış ve gerçek bir web sistemi üzerinde analiz edilmiştir. Geliştirilen Açgözlü atama algoritması, veri boyutu eksponansiyel olarak arttığında bile sistemin çökmesini engellemiş ve performansı korumuştur. Sonuç olarak, bu tür gerçek dünya NP-Zor problemlerinde kesin çözümler yerine optimuma yakın (near-optimal) sonuç veren sezgisel yaklaşımların kullanımı, sistem kullanılabilirliği açısından zorunludur.

\begin{thebibliography}{00}
\bibitem{b1} E. K. Burke, J. P. Newall, and R. F. Weare, ``A memetic algorithm for university exam timetabling,'' \textit{Practice and Theory of Automated Timetabling}, vol. 1153, pp. 241-250, 1996.
\bibitem{b2} M. R. Garey and D. S. Johnson, \textit{Computers and Intractability: A Guide to the Theory of NP-Completeness}. W. H. Freeman, 1979.
\end{thebibliography}

\end{document}
